\subsection{谱的孤立点}

设 E 为 Banach 空间，u 为 E 的自同态。设 \(\lambda\in C\) 为 \(\mathrm{Sp}(u)\) 的孤立点。我们记 \(\mathrm{E}_{\lambda}(u)=\mathrm{E}_{\{\lambda\}}(u)\) ，并用 \(e_{\lambda}(u)\) 表示以 \(\mathrm{E}_{\lambda}(u)\) 为像的、与 u 和 \(\{\lambda\}\) 相伴的谱投影。

我们还记 \(\widetilde{\mathrm{E}}_{\lambda}(u) = \mathrm{E}_{\mathrm{Sp}(u) - \{\lambda\}}(u)\) 。\(\widetilde{\mathrm{E}}_{\lambda}(u)\) 的由 \(u\) 诱导的自同态的谱是 \(\mathrm{Sp}(u) - \{\lambda\}\) ；特别地，\(u - \lambda 1_{\mathrm{E}}\) 诱导出 \(\widetilde{\mathrm{E}}_{\lambda}(u)\) 的一个自同构。

空间 \(\mathrm{E}_{\lambda}(u)\) 非零。由 u 诱导的 \(\mathrm{E}_{\lambda}(u)\) 的自同态的谱退化为 \(\lambda\) ，因此 \(u - \lambda1_{E}\) 诱导出 \(\mathrm{E}_{\lambda}(u)\) 的一个拟幂零自同态。自同态 \(u - \lambda1_{E}\) 诱导出 \(\widetilde{\mathrm{E}}_{\lambda}(u)\) 的一个自同构，因此 \(\mathrm{Ker}(u - \lambda1_{\mathrm{E}})^{n} \subset \mathrm{E}_{\lambda}(u)\) 对所有 \(n \in N\) 成立。特别地，我们有 \(\mathrm{Ker}(u - \lambda1_{\mathrm{E}}) \subset \mathrm{E}_{\lambda}(u)\) 。

\(\lambda\) 为阶 \(p > 0\) 的极点，即 \(u\) 的预解式的极点，其充分必要条件是 \((u - \lambda 1_{\mathrm{E}})^{p - 1}e_{\lambda}(u)\neq 0\) 且 \((u - \lambda 1_{\mathrm{E}})^{p}e_{\lambda}(u) = 0\) （I，第 83 页，命题 17 的推论）。这时，我们有 \(\mathrm{E}_{\lambda}(u) = \mathrm{Ker}((u - \lambda 1_{\mathrm{E}})^{p})\) 与 \(\widetilde{\mathrm{E}}_{\lambda}(u) = \mathrm{Im}((u - \lambda 1_{\mathrm{E}})^{p})\) ，因为 \((u - \lambda 1_{\mathrm{E}})^{p}\) 诱导出 \(\widetilde{\mathrm{E}}_{\lambda}(u)\) 的一个自同构。我们还有

\[
(u - \lambda 1 _ {\mathrm{E}}) ^ {p - 1} e _ {\lambda} (u) = \lim _ {z \to \lambda} (z - \lambda 1 _ {\mathrm{E}}) ^ {p} \mathrm{R} (u, z)
\]

这出于 I，第 83 页，命题 17。

\pandocbounded{\includegraphics[keepaspectratio]{images/part001_eb2e85aa.jpg}}

注意，一般说来 \(\mathrm{E}_{\lambda}(u)\) 并不是族 \((\mathrm{Ker}((u-\lambda1_{\mathrm{E}})^{n}))_{n\in\mathbf{N}}\) 的并，甚至也不是这个并的闭包；特别地，\(\mathrm{Sp}(u)\) 的孤立点不一定是 u 的特征值（I，第 187 页，习题 1）。同样，也可能存在 u 的不是 \(\mathrm{Sp}(u)\) 的孤立点的特征值（I，第 188 页，习题 2）。

